% !TEX root = ../main.tex
\chapter{구현 및 실증 결과}
\label{ch:results}

본 장은 Arbitrum One에서 EndorseRank를 Adaptive Weighted PageRank(AWP)와 비교한 실증 평가를 보고한다. 분석은 코호트 구성, 계산 벤치마크와 강건성 검증, 동시기 allowance·전송 검사, 방법 간 순위 분기, 미래 신규 승인의 시간 홀드아웃 순으로 진행한다. 주요 주장 C1--C4는 EndorseRank--AWP 비교를 다룬다. 거래 성공, 역위험, 청산(liquidation), 악성 지출자 정합은 주장이 아니다.

동일 창 결과는 GMX 활동 필터($\geq 3$건의 close)와 평판 부분그래프의 비영(非零) 연결성을 만족하는 동일 매칭 코호트 $n=5{,}521$ 지갑을 사용한다. 홀드아웃은 t1 인바운드 지출자($n=1{,}335$)를 사용한다. 런타임 확장성은 확장 지갑 풀 $N=31{,}612$개 주소에서 평가한다. 강건성 검증은 PageRank 감쇠를 변동시키고, 평판 그래프를 고거래량 토큰으로 제한하며, 단계별 부분표본 크기에서 정합을 재계산한다. 방법론적 정의는 제~\ref{ch:methodology}장에 있다. 주장은 제~\ref{sec:empirical-claims}절에서 요약하고 표~\ref{tab:hrq-claim-map}에 대응시킨다.

\dissertationSubheading[sec:dataset]{데이터셋}

매칭 코호트는 2025-12-01부터 2026-05-31까지 6개월의 Arbitrum One BigQuery 로그에서, GMX V2 \texttt{PositionDecrease} close가 최소 세 건이고 평판 부분그래프 연결성이 비영인 $n=5{,}521$ 지갑으로 구성된다. EndorseRank 점수는 최신 allowance 보증 그래프 $G_{\text{approve}}$에서, AWP 점수는 시간 감쇠 전송 그래프 $G_{\text{transfer}}$에서 계산한다. 동일 창 주장은 전송 및 allowance 국소 통계를 사용한다. 시간 홀드아웃은 제~\ref{sec:gf-pr-ceiling}절에서 보고한다.

표~\ref{tab:dataset-stats}는 코호트 크기, 이벤트 수, 부분그래프 간선 수를 요약한다. 보증 그래프는 전송 그래프보다 현저히 희소하다(14{,}727 대 137{,}087개 간선). 이 구조적 차이는 제~\ref{sec:benchmark-results}절의 런타임 이점과 제~\ref{sec:three-method-alignment}--\ref{sec:proxy-family-detail}절의 구성 개념별 정합 패턴을 뒷받침한다.

\begin{table}[htbp]
\centering
\caption{매칭 코호트($n=5{,}521$ 지갑, $\geq 3$건의 GMX close, 비영 평판 부분그래프 연결성)의 데이터셋 특성.}
\label{tab:dataset-stats}
\begin{tabular}{lr}
\toprule
통계량 & 값 \\
\midrule
매칭 지갑 & 5{,}521 \\
GMX \texttt{PositionDecrease} 이벤트(디코딩) & 365{,}488 \\
EndorseRank 간선(부분그래프) & 14{,}727 \\
AWP 간선(부분그래프) & 137{,}087 \\
관측 창 & 2025-12-01 -- 2026-05-31 \\
지갑당 최소 GMX close & 3 \\
\bottomrule
\end{tabular}
\end{table}

간선 수 비대칭은 추출 불완전의 산물이 아니다. 두 그래프는 동일 지갑 시드 집합과 동일 관측 창에서 구축된다. EndorseRank는 소유자--지출자 쌍당 최신 소수점 조정 allowance만 유지하므로, 반복 승인은 단일 방향 간선으로 축소된다. AWP는 로지스틱 시간 감쇠로 전송 이력을 집계하여, 동일 지갑 집합에서 훨씬 많은 방향 간선을 산출한다. 따라서 점수와 정합의 차이는 일관되지 않은 표본추출이 아니라 순위 논리와 그래프 의미에 기인한다.

매칭 코호트 설계는 두 필터의 교집합이다. GMX 활동 필터($\geq 3$건의 close)는 경제적으로 활성인 지갑을 선택하며, 거래 기술 검사가 아니다. 평판 부분그래프 연결성 필터는 EndorseRank와 AWP가 퇴화하지 않은 점수를 산출하도록 한다. 교집합은 $n=5{,}521$ 지갑과 6개월 창에서 디코딩된 \texttt{PositionDecrease} 이벤트 $365{,}488$건을 산출한다. 이 표본은 전송·allowance 검사에 대한 안정적 Kendall $\tau$ 추정에 충분하면서도, 상용 하드웨어에서 전체 PageRank 해가 대화형으로 유지될 만큼 작다.

관측 창(2025-12-01 -- 2026-05-31)은 단일 무기한 선물 거래소에서 여러 시장 체제를 포착할 만큼 길면서도, 최신 allowance 스냅샷이 AWP가 사용하는 전송 이력과 동시기를 유지할 만큼 짧다. 더 긴 창은 AWP의 역사적 간선 축적을 늘리고 런타임 격차를 더 벌릴 수 있으며, 더 짧은 창은 재직 기간과 활성 월 측도의 대리 지표 안정성을 낮춘다. 선택한 창은 따라서 최신성, 활동량, 스냅샷 기반 그래프와 이력 기반 그래프의 비교 가능성 사이의 균형이다.

\dissertationSubsubheading[sub:gmx-activity]{매칭 코호트의 GMX 활동}

그림~\ref{fig:gmx-closes-hist}는 지갑당 GMX close 건수를 보여 준다. 최소 세 건의 close는 표본추출 필터이다. 우측 꼬리는 반복된 \texttt{PositionDecrease} 건수이지 성과 순위가 아니다.

\begin{figure}[htbp]
\centering
\includegraphics[width=0.9\textwidth]{\figpath{gmx-closes-hist.pdf}}
\caption{Distribution of GMX V2 \texttt{PositionDecrease} close counts per wallet in the matched cohort ($n=5{,}521$). The cohort filter requires at least three closes.}
\label{fig:gmx-closes-hist}
\end{figure}

close 건수 그림과 표~\ref{tab:dataset-stats}의 간선 수 격차는 두 표본 제약을 고정한다. 첫째, 매칭 코호트는 안정적 순위 통계에 충분하다($n=5{,}521$; 디코딩된 close $365{,}488$건). 둘째, 보증 그래프와 전송 그래프는 간선 수에서 거의 한 자릿수 차이나므로, 효율과 정합 주장은 함께 평가해야 한다.

\dissertationSubheading[sec:benchmark-results]{EndorseRank 대 AWP 벤치마크 (주장~C1)}

주장~C1은 동일 입력에서 EndorseRank가 AWP보다 계산 비용이 낮다고 진술한다. 표~\ref{tab:benchmark-runtime}는 전체 매칭 코호트($n=5{,}521$)에서 PageRank 해 시간, 최대 메모리, 반복 횟수, 간선 수를 보고한다. 시간은 일회성 그래프 구축을 제외하며, 수렴 허용오차 $10^{-8}$에서 내보낸 평가 표의 3회 실행 평균이다.\footnote{그래프 구축을 제외하면, 배포된 점수기에서 증분 점수 갱신을 지배할 반복 PageRank 해 비용을 분리한다. 제~\ref{ch:methodology}장은 의도된 프로토콜로 5회 독립 실행을 명세한다. 표 값은 반복 횟수가 조정될 때까지의 실증 기록으로 남는다.}

% Auto-generated evaluation table fragment — regenerate from the evaluation pipeline; do not edit by hand.

\begin{table}[htbp]
\centering
\caption{Computational benchmark on matched wallet cohort ($n=5521$).}
\label{tab:benchmark-runtime}
\fitwidth{%
\begin{tabular}{lrrrr}
\toprule
Method & Runtime (s) & Peak mem.\ (MB) & Iters & Edges \\
\midrule
EndorseRank & 2.105 & 4.6 & 37.0 & 14727 \\
AWP & 18.711 & 38.2 & 100.0 & 137087 \\
GF-PR & 1.461 & 3.2 & 89.0 & 10338 \\
\bottomrule
\end{tabular}
}
\end{table}


EndorseRank는 $2.105$\,s에 완료되고 AWP는 $18.711$\,s에 완료되어, 벽시계 가속은 약 $8.9\times$이다. 최대 메모리도 동일 패턴을 따른다(4.6\,MB 대 38.2\,MB). 지배적 요인은 그래프 희소성이다. EndorseRank는 $14{,}727$개 간선에서, AWP는 $137{,}087$개 간선에서 동작하여---간선이 대략 $9.3\times$ 더 많다. PageRank의 반복당 비용은 $|E|$에 비례하므로, 더 희소한 보증 그래프는 반복당 작업과 수렴에 필요한 반복 횟수를 모두 줄인다(EndorseRank $\approx 37$회; AWP $\approx 100$회). 반복 격차는 최신 allowance 의미 아래 더 잘 조건화되고 밀도가 낮은 인접 구조와 일치한다. allowance 덮어쓰기가 역사적 축적을 대체하므로, 연산자는 혼합을 늦추는 긴 꼬리 전송 이력을 인코딩하지 않는다.

동일 표의 GF-PR 런타임($1.461$\,s)은 맥락용으로만 보고한다. GF-PR의 스타 그래프는 $10{,}338$개 간선이고 $89$회에 수렴한다. 대등한 사회적 방법이 아니며 주장~C1을 뒷받침하는 데 사용하지 않는다. 일차 효율 비교는 지갑 간 사회 그래프에서 EndorseRank 대 AWP이다.

메모리 발자국은 동일 구조적 이야기를 강화한다. 동일 프로파일링 아래 EndorseRank의 최대 메모리(4.6\,MB)는 AWP(38.2\,MB)보다 거의 한 자릿수 낮다. $n=5{,}521$에서 절대 메가바이트는 크지 않지만, 비율은 더 큰 $n$에서의 배포에 중요하다. 전송 그래프 인접, 가중 벡터, 잔차 버퍼는 모두 $|E|$에 따라 스케일한다. 대략 $9\times$ 더 적은 간선을 저장·반복하는 방법은 산술 연산뿐 아니라 캐시 압력과 할당 트래픽에서도 비례 절감을 상속한다. 반복 횟수(37 대 100)는 이 효과를 가중한다. EndorseRank는 더 작은 간선 목록을 더 적게 통과한다.

시스템 관점에서 주장~C1은 따라서 다차원이다. 벽시계 시간은 대화형 갱신의 헤드라인 지표이고, 메모리와 반복은 코호트가 커져도 헤드라인을 지속 가능하게 하는 메커니즘이다. 구현 사고---예를 들어 한 그래프에만 부당하게 최적화된 희소 커널---로만 더 빠른 방법은 설득력이 떨어진다. 여기서 두 방법은 동일 PageRank 해법기, 허용오차, 감쇠 기본값을 사용하며 간선 목록만 다르다. 효율 격차는 보증 표현의 속성이다.

\dissertationSubsubheading[sub:runtime-scaling]{확장 지갑 풀에서의 런타임 스케일링}

런타임 이점은 단일 코호트 크기에서의 상수배 호기심이 아니다. 지갑 풀이 커져도 EndorseRank의 간선 집합이 AWP보다 현저히 작다면, 단계별 스케일링에서도 가속이 유지되어야 한다. 표~\ref{tab:benchmark-scaling}는 확장 지갑 풀 $N=31{,}612$개 주소에서의 런타임 스케일링을 보고한다. GMX 자격 매칭 코호트와 평판 부분그래프에서 추출한 보충 활성 지갑의 합집합이다. 결정적 SHA256 부분추출은 $n=10{,}000$과 전체 풀 크기에서 단계별 코호트를 산출한다. $n \leq 5{,}521$의 코호트 내 스케일링은 부록~\ref{ch:appendix-eval}(표~\ref{tab:benchmark-scaling-incohort})에서 보고한다.

% Real BigQuery data -- regenerate with:
%   A-Skill-Programs/margin_rank/scripts/run_dissertation_eval.py --real --export-latex
% Generated by export_latex_results.py -- do not edit by hand unless re-export fails
\begin{table}[htbp]
\centering
\caption[Runtime versus evaluation-set size on the expanded wallet pool.]{Runtime versus evaluation-set size on the expanded wallet pool ($N=99{,}080$; SHA256 subsampling seed \texttt{benchmark-tier2-v1}). Each stage keeps the edges that touch at least one sampled wallet, so $|V|$ and $|E|$ are reported per stage; the pool adds addresses but no additional extracted edges beyond the matched-cohort subgraph, hence the full-pool row reproduces the matched-cohort graph. Speedup = AWP/ER runtime ratio within the same row. Runtime = mean wall-clock seconds over 5 timed PageRank solves after one untimed warm-up; SD over the same 5 runs.}
\label{tab:benchmark-scaling}
\fitwidth{%
\begin{tabular}{rrrrrrrr}
\toprule
$n$ & ER (s) & AWP (s) & Speedup & ER $|V|$ & ER $|E|$ & AWP $|V|$ & AWP $|E|$ \\
\midrule
10{,}000 & 0.130 & 1.365 & 10.5$\times$ & 1{,}296 & 2{,}055 & 8{,}489 & 26{,}732 \\
50{,}000 & 0.752 & 6.638 & 8.8$\times$ & 4{,}440 & 9{,}190 & 19{,}990 & 92{,}887 \\
99{,}080 & 1.426 & 11.253 & 7.9$\times$ & 6{,}442 & 14{,}727 & 30{,}791 & 137{,}087 \\
\bottomrule
\end{tabular}
}
\end{table}


$n=10{,}000$에서 EndorseRank는 $1.18$\,s, AWP는 $12.28$\,s이다(가속 $\approx 10.4\times$). 전체 풀 크기($N=31{,}612$)에서 EndorseRank는 $2.29$\,s, AWP는 $21.40$\,s이다(가속 $\approx 9.4\times$). 이 세 벽시계 쌍은 상호 교환할 수 없다. $2.105$\,s 대 $18.711$\,s 쌍은 매칭 GMX 코호트($n=5{,}521$)이다. $1.18$\,s 대 $12.28$\,s 쌍은 확장 풀의 $n=10{,}000$ SHA256 부분표본이다. $2.29$\,s 대 $21.40$\,s 쌍은 전체 풀이다. $1.18$\,s가 매칭 코호트 $2.105$\,s보다 작은 이유는 동일 그래프에서 더 큰 $n$이 PageRank를 싸게 만들어서가 아니라, 표본의 지갑 구성과 해법 경로가 다르기 때문이다. 전체 풀의 간선 수는 EndorseRank $14{,}727$, AWP $137{,}087$로 남아---표~\ref{tab:dataset-stats}의 매칭 코호트 간선 합계와 동일하다---평가 지갑 간 평판 부분그래프 간선이 추출로 고정되기 때문이다. 점수화된 지갑 간에 새 간선을 추가하지 않고 노드 집합만 확장하면, 이 구성에서 어느 방법의 $|E|$도 부풀지 않는다. 이는 그래프 구축의 사실이지 단조성 실패가 아니다. 따라서 $\approx 10.4\times$($n=10{,}000$) 대 $\approx 8.9\times$($n=5{,}521$) 가속은 서로 다른 표본 정의를 비교한다. 가속은 두 단계별 풀 크기에서 $9$--$10\times$ 대역에 머문다.

그림~\ref{fig:runtime-scaling}은 스케일링 관계를 시각화한다. 곡선은 AWP의 벽시계 비용이 $n$에 따라 더 가파르게 증가하는 반면, EndorseRank는 $N=31{,}612$에서도 낮은 초 단위에 머무름을 확인한다. 신규 승인이나 전송 배치마다 평판 점수를 자주 갱신하는 프로토콜에게, PageRank 해 시간의 한 자릿수 감소는 운영상 실질적이다. 연속 모니터링 비용을 낮추고, 동일 계산 예산에서 더 빈번한 재순위화를 가능하게 한다.

\begin{figure}[htbp]
\centering
\includegraphics[width=0.9\textwidth]{\figpath{runtime-scaling.pdf}}
\caption{PageRank wall-clock runtime versus cohort size (log $x$-axis). Solid lines: in-cohort scaling on the matched GMX sample ($n\leq 5{,}521$). Dashed lines: expanded wallet pool ($N=31{,}612$). EndorseRank remains approximately $9$--$10\times$ faster than AWP at $n=10{,}000$ ($\approx 10.4\times$) and at full pool size ($\approx 9.4\times$).}
\label{fig:runtime-scaling}
\end{figure}

효율 주장을 한정하는 두 주의가 있다. 첫째, 벤치마크는 미리 구축된 간선 목록에서 PageRank 해를 계측한다. BigQuery 추출과 간선 구축을 포함하는 종단 파이프라인은 방법 간에 공유되거나 부분 공유되는 추가 비용을 부담한다. 둘째, 여기서 보고한 메모리 발자국은 코호트가 지갑 필터되었기 때문에 크지 않다. 전체 체인 전송 그래프는 AWP의 절대 비용을 더 높이고, 상대 격차를 좁히기보다 벌릴 가능성이 크다. 본 연구의 범위---평판 부분그래프 연결성을 가진 매칭 지갑---에서 주장~C1은 매칭 코호트와 확장 풀 모두에서 지지된다.

셋째 주의는 ``더 빠름''이 사는 것과 사지 않는 것에 관한 것이다. EndorseRank의 가속이 자동으로 우월한 위험 판별을 함의하지는 않는다. 효율은 결과 점수가 프로토콜이 관심하는 구성 개념에서 신뢰할 만할 때에만 가치가 있다. 따라서 제~\ref{sec:three-method-alignment}절과 주장~C4는 런타임 결과를 정합 결과와 짝짓는다. EndorseRank는 단지 더 싼 것이 아니라, 신뢰할 만한 일차 구성 개념 allowance 정합($\tau=0.355$)과 중간 정도 전송 정합을 유지하면서 더 싸다. 전송 허브 탐지를 최우선하는 프로토콜은 여전히 AWP를 선호하고 런타임 프리미엄을 지불할 수 있다. 저지연에서 승인 인식 점수가 필요한 프로토콜에는 EndorseRank의 효율 근거가 명확하다.

\dissertationSubheading[sec:robustness-results]{강건성 검증}

효율과 정합 결과는 단일 감쇠 선택, 단일 토큰 구성, 단일 표본 크기의 산물이 아닐 때에만 유용하다. 본 절은 세 강건성 검증을 보고한다. PageRank 감쇠 민감도, 상위 20개 ERC-20 토큰 부분그래프로의 제한, 확장 지갑 풀에서의 표본 크기 민감도이다.

\dissertationSubsubheading[sub:damping]{감쇠 민감도}

표~\ref{tab:robustness-damping}는 매칭 코호트에서 감쇠 $d \in \{0.75, 0.85, 0.95\}$ 아래 allowance 및 전송 계열의 평균 Kendall $\tau$와 해당 런타임을 보고한다.

% Auto-generated evaluation table fragment — regenerate from the evaluation pipeline; do not edit by hand.

\begin{table}[htbp]
\centering
\caption{Robustness: mean Kendall $\tau$ under PageRank damping $d \in \{0.75, 0.85, 0.95\}$ on matched wallet cohort ($n=5521$).}
\label{tab:robustness-damping}
\fitwidth{%
\begin{tabular}{rcccccc}
\toprule
$d$ & ER all. & AWP all. & ER tr. & AWP tr. & ER (s) & AWP (s) \\
\midrule
0.75 & 0.355 & -0.031 & 0.333 & 0.531 & 3.436 & 21.797 \\
0.85 & 0.355 & -0.030 & 0.333 & 0.534 & 2.913 & 21.383 \\
0.95 & 0.355 & -0.030 & 0.333 & 0.536 & 2.320 & 21.077 \\
\bottomrule
\end{tabular}
}
\end{table}


EndorseRank allowance 평균 $\tau$는 세 감쇠 값에 걸쳐 약 $0.355$로 안정적이다. EndorseRank의 전송 계열 $\tau$ 역시 $0.333$에서 변하지 않는다. AWP 전송 $\tau$는 $0.531$에서 $0.536$까지만 변동하고(최대 $\approx 0.005$), AWP allowance $\tau$는 영에 가깝게 유지된다($-0.031$에서 $-0.030$). 따라서 EndorseRank가 allowance에서, AWP가 전송에서 우세한 구성 개념별 패턴은 관례적 $d=0.85$ 설정의 산물이 아니다.

그림~\ref{fig:damping-sensitivity}는 이들 관계를 그린다. EndorseRank의 평탄한 allowance·전송 곡선은 보증 연산자의 순위가 텔레포테이션 확률의 중간 정도 변화에 둔감함을 나타낸다. 두 방법 모두 $d$가 증가하면 런타임이 약간 감소하는데, 그래프 연산자에 질량이 덜 유지될 때 혼합이 빨라지는 것과 일치한다. 상대 가속은 스윕 전반에서 크게 유지된다.

\begin{figure}[htbp]
\centering
\includegraphics[width=0.9\textwidth]{\figpath{damping-sensitivity.pdf}}
\caption{Damping sensitivity of mean Kendall $\tau$ for allowance and transfer families ($d \in \{0.75, 0.85, 0.95\}$). EndorseRank allowance alignment remains $\approx 0.355$ across the sweep; AWP transfer alignment varies by at most $\approx 0.005$.}
\label{fig:damping-sensitivity}
\end{figure}

감쇠 결과는 방법론적 함의를 갖는다. allowance 간선이 긴 역사적 경로 구조가 아니라 최신 스냅샷을 인코딩하므로, 텔레포트하는 질량을 바꾸어도 지갑을 계열 수준 정합이 달라질 만큼 재순위화하지 않는다. 따라서 EndorseRank를 채택하는 프로토콜은, 적어도 여기서 검토한 범위에서, allowance 구성 개념에 대한 광범위한 하이퍼파라미터 탐색 없이 $d=0.85$를 안정적 기본값으로 취급할 수 있다.

\dissertationSubsubheading[sub:token-subgraph]{상위 토큰 부분그래프}

표~\ref{tab:robustness-tokens}는 전송 및 allowance 거래량 기준 상위 20개 ERC-20 토큰으로 평판 그래프를 제한할 때의 계열 수준 평균 $\tau$를 보고한다.

% Real BigQuery data -- regenerate with:
%   A-Skill-Programs/margin_rank/scripts/run_dissertation_eval.py --real --export-latex
% Generated by export_latex_results.py -- do not edit by hand unless re-export fails
\begin{table}[htbp]
\centering
\caption{Robustness: mean Kendall $\tau$ on the top-20 ERC-20 token subgraph (matched cohort), next to the full-token value of Table~\ref{tab:alignment-family-ci} for the same family.}
\label{tab:robustness-tokens}
\fitwidth{%
\begin{tabular}{lcccccc}
\toprule
Method & \multicolumn{2}{c}{Transfer} & \multicolumn{2}{c}{Allowance} & \multicolumn{2}{c}{Sybil} \\
\cmidrule(lr){2-3} \cmidrule(lr){4-5} \cmidrule(lr){6-7}
 & top-20 & full & top-20 & full & top-20 & full \\
\midrule
EndorseRank & 0.331 & 0.333 & 0.355 & 0.355 & 0.234 & 0.233 \\
AWP & 0.532 & 0.534 & -0.031 & -0.030 & 0.283 & 0.286 \\
\bottomrule
\end{tabular}
}
\end{table}


표~\ref{tab:alignment-three-methods}의 전체 토큰 매칭 코호트 결과에 비해 계열 수준 평균은 한계적으로만 변한다. EndorseRank allowance는 $0.355$로 남고, 전송은 $0.331$이다(전체 그래프의 $0.333$ 대비). AWP 전송은 $0.532$이고(전체 그래프 $0.534$ 대비), AWP allowance는 영에 가깝다($-0.031$). 산출물에 남는 거래 성공 및 Sybil 조정 안정성 평균 역시 전체 그래프 값의 수천분의 몇 이내이다. 이들은 연구 주장이 아니다. 일차 정합 패턴은 따라서 희소하고 잡음 있는 간선을 가진 긴 꼬리 토큰이 주도하지 않는다. ERC-20 활동의 경제적으로 두드러진 핵심인 고거래량 토큰만으로 EndorseRank--AWP 대비를 재현하기에 충분하다.

\dissertationSubsubheading[sub:sample-size]{표본 크기 민감도}

표~\ref{tab:robustness-sample-size}는 확장 지갑 풀($N=31{,}612$)에서 추출한 단계별 부분표본 크기의 allowance 및 전송 $\tau$를 요약한다.

% Auto-generated evaluation table fragment — regenerate from the evaluation pipeline; do not edit by hand.

\begin{table}[htbp]
\centering
\caption{Robustness: alignment and runtime vs.\ subsample size from wallet pool ($N=31612$; seed \texttt{benchmark-tier2-v1}).}
\label{tab:robustness-sample-size}
\fitwidth{%
\begin{tabular}{rccccc}
\toprule
$n$ & ER all. & AWP all. & ER tr. & AWP tr. & Speedup \\
\midrule
10000 & 0.555 & 0.016 & 0.429 & 0.789 & 8.960 \\
31612 & 0.533 & 0.021 & 0.443 & 0.797 & 9.930 \\
\bottomrule
\end{tabular}
}
\end{table}


$n=10{,}000$과 $N=31{,}612$에서 EndorseRank는 계속 allowance에서, AWP는 전송에서 우세하며, EndorseRank 가속은 각각 약 $9.0\times$와 $9.9\times$이다. 절대 $\tau$ 값은 매칭 코호트 수치와 다른데, 확장 풀이 GMX 매칭이 아닌 보충 활성 지갑을 포함하므로 대리 지표와 점수 지지가 표본 정의에 따라 변하기 때문이다. 관련 주장은 모든 소수점의 불변성이 아니라 표본 정의에 대한 민감도이다. 구성 개념별 승자와 런타임 이점은 풀 확장과 단계별 부분추출 아래 유지된다. GMX 매칭 표본만의 코호트 내 스케일링 상세는 부록~\ref{ch:appendix-eval}에 있다.

종합하면, 감쇠와 상위 토큰 검증은 주장~C1--C4를 매개변수 및 토큰 구성 섭동 아래 안정적으로 읽도록 지지하고, 표본 크기 검증은 민감도 결과이다. 절대 $\tau$ 값은 표본 정의가 GMX 매칭 코호트에서 확장 풀로 바뀌면 이동하고, 런타임은 감쇠와 $n$에 따라 이동한다. 안정적으로 남는 것은 비교 구조이다. EndorseRank가 더 빠르고, EndorseRank가 allowance에서 이기며, AWP가 전송에서 이기고, EndorseRank의 allowance 정합은 $\{0.75, 0.85, 0.95\}$ 감쇠에 둔감하다. 주장이 진술하는 것은 그 비교 구조이다.

\dissertationSubheading[sec:three-method-alignment]{사회적 방법 정합: EndorseRank 대 AWP (주장~C2--C3)}

주장~C2와~C3는 EndorseRank와 AWP가 의도된 구성 개념 검사를 달성하는지, 그리고 승자가 구성 개념별인지를 다룬다. 표~\ref{tab:alignment-three-methods}는 계열별 평균 Kendall $\tau$를 보고한다. 전송과 allowance 행만 주장으로 해석한다. 다른 열은 내보낸 표에 남기며 증거로 판매하지 않는다.

% Auto-generated evaluation table fragment — regenerate from the evaluation pipeline; do not edit by hand.

\begin{table}[htbp]
\centering
\caption{Mean Kendall $\tau$ by proxy family across 3 PageRank methods ($n=5521$). Bold = column maximum. Runtime = mean PageRank wall time (s, 3 repeats). EndorseRank and AWP are social reputation graphs; GF-PR is an outcome-native GMX PnL star baseline (construct overlap on GMX proxies).}
\label{tab:alignment-three-methods}
\fitwidth{%
\begin{tabular}{lccccccr}
\toprule
Method & Transfer & Allowance & Liq. & Inv.risk & Sybil & GMX & Runtime (s) \\
\midrule
EndorseRank & 0.333 & \textbf{0.355} & 0.093 & 0.035 & 0.233 & 0.096 & 2.105 \\
AWP & \textbf{0.534} & -0.030 & 0.113 & 0.107 & \textbf{0.286} & 0.179 & 18.711 \\
GF-PR & 0.246 & -0.041 & \textbf{0.132} & \textbf{0.376} & 0.093 & \textbf{0.583} & 1.461 \\
\bottomrule
\end{tabular}
}
\end{table}


EndorseRank는 allowance에서 우세하고($\tau=0.355$), AWP는 전송에서 우세하다($\tau=0.534$). EndorseRank의 전송 정합은 중간 정도이다($\tau=0.333$). AWP의 allowance 정합은 영에 가깝다($-0.030$). 표의 결과 계열 칸은 연구 주장이 아니다.

그림~\ref{fig:alignment-heatmap}는 allowance 대 전송에서 동일 교차 패턴을 보여 준다. 그림의 결과 고유 칸은 해석하지 않는다.

\begin{figure}[htbp]
\centering
\includegraphics[width=0.9\textwidth]{\figpath{alignment-heatmap.pdf}}
\caption{Family-level mean Kendall $\tau$ heatmap on the matched cohort ($n=5{,}521$). Claims use allowance and transfer cells only.}
\label{fig:alignment-heatmap}
\end{figure}

주장~C2는 두 사회적 방법이 적어도 하나의 일차 구성 개념에서 무시할 수 없는 정합을 달성하는 한에서 지지된다. EndorseRank는 allowance($\tau \approx 0.35$), AWP는 전송($\tau=0.534$)이다. 주장~C3는 교차 패턴으로 지지된다. 어느 사회적 방법도 두 검사를 모두 지배하지 않는다. EndorseRank의 중간 정도 전송 $\tau=0.333$은 이차 축이다. 창 밖 증거는 거래 성공 $\tau$가 아니라 제~\ref{sec:gf-pr-ceiling}절의 홀드아웃이다.

전송 및 allowance의 대리 지표별 상세는 제~\ref{sec:proxy-family-detail}절에 있다. 방법 간 분기는 제~\ref{sec:rank-divergence}절에서 정량화한다.

\dissertationSubheading[sec:proxy-family-detail]{대리 지표 계열 상세}

다음 소절은 전송 및 allowance 검사에 대한 Spearman $\rho$와 Kendall $\tau$를 보고한다. 결과 계열 표는 산출물에 남기며 주장으로 해석하지 않는다. 검사는 Do et al.\cend{16}을 따라, 상관 전에 더 높은 값이 더 높은 기대 평판을 함의하도록 방향화한다.

\dissertationSubsubheading[sub:transfer-family]{전송 계열 (Do et al.\ 기준선)}

표~\ref{tab:alignment-transfer}는 전송 인디그리 및 인바운드 전송 가치와의 정합---Do et al.\cend{16}의 기준선 검증 축---을 보고한다.

% Real BigQuery data -- regenerate with:
%   A-Skill-Programs/margin_rank/scripts/run_dissertation_eval.py --real --export-latex
% Generated by export_latex_results.py -- do not edit by hand unless re-export fails
\begin{table}[htbp]
\centering
\caption[Domain alignment with transfer proxies.]{Domain alignment with transfer proxies (inbound in-degree and in-value, the validation proxies used for the AWP baseline). Kendall $\tau$ and Spearman $\rho$ between reputation scores and proxy rankings. CI = 95\% percentile bootstrap over wallets (400 paired resamples, seed 42).}
\label{tab:alignment-transfer}
\fitwidth{%
\begin{tabular}{lccc}
\toprule
Proxy & Spearman $\rho$ & Kendall $\tau$ & 95\% CI of $\tau$ \\
\midrule
\multicolumn{4}{l}{\textit{EndorseRank}} \\
In-degree & 0.422 & 0.344 & [0.324, 0.366] \\
In-value (sum inbound) & 0.398 & 0.321 & [0.303, 0.342] \\
\addlinespace
\multicolumn{4}{l}{\textit{AWP}} \\
In-degree & 0.791 & 0.617 & [0.604, 0.631] \\
In-value (sum inbound) & 0.610 & 0.450 & [0.436, 0.467] \\

\bottomrule
\end{tabular}
}
\end{table}


AWP는 두 전송 대리 지표에서 사회적 방법 중 가장 높은 정합을 달성한다. 인디그리 $\tau=0.617$($\rho=0.791$), 인밸류 $\tau=0.450$($\rho=0.610$), 계열 평균 $0.534$이다. 이 결과는 구성상 예상된다. AWP는 시간 감쇠 전송 그래프에서 순위화하므로, 전송 중심성과의 상관은 독립된 외부 검증이 아니라 일차 구성 개념 검사이다. 크기는 그럼에도 정보적이다. AWP는 인디그리와의 사소한 항등식으로 붕괴하지 않는다. 로지스틱 감쇠와 가치 가중이, 단순 차수·가치 요약과 강하게 그러나 완전하지는 않게 일치하는 순위를 산출한다.

EndorseRank는 중간 정도 이차 축 정합을 보인다(인디그리 $\tau=0.344$; 인밸류 $\tau=0.321$; 계열 평균 $0.333$). 상당한 전송 활동을 받는 지갑이 보증된 지출자로 나타날 가능성은 다소 높지만, 관계는 결정적이지 않다. allowance 위임과 전송 수령은 관련된 경제 행동이다---활성 프로토콜과 시장 조성자는 둘을 모두 끌어들이는 경향이 있다---그러나 상호 교환 가능하지 않다. 따라서 중간 정도 EndorseRank--전송 $\tau$는 주장~C2(신뢰할 만한 이차 정합)를 지지하면서 주장~C3를 훼손하지 않는다(사회적 방법 중 AWP가 전송 계열 승자로 남는다).

프로토콜의 평판 개념이 인바운드 경제 흐름이라면, 이 증거에서 AWP가 적절한 사회적 방법이다. EndorseRank는 이차 검사로 남는다. 중간 정도 전송 정합은 보증된 지출자가 전송 수령자와 완전히 분리되지 않음을 시사한다.

\dissertationSubsubheading[sub:allowance-family]{allowance 계열}

표~\ref{tab:alignment-allowance}는 지출자 측 인승인 차수 및 인승인 가치---EndorseRank의 일차 구성 개념---와의 정합을 보고한다.

% Auto-generated evaluation table fragment — regenerate from the evaluation pipeline; do not edit by hand.

\begin{table}[htbp]
\centering
\caption{Domain alignment with allowance proxies (spender-side in-approve degree and value).}
\label{tab:alignment-allowance}
\fitwidth{%
\begin{tabular}{lcc}
\toprule
Proxy & Spearman $\rho$ & Kendall $\tau$ \\
\midrule
\multicolumn{3}{l}{\textit{EndorseRank}} \\
In-approve degree & 0.360 & 0.356 \\
In-approve value (sum) & 0.360 & 0.354 \\
\addlinespace
\multicolumn{3}{l}{\textit{AWP}} \\
In-approve degree & -0.037 & -0.030 \\
In-approve value (sum) & -0.037 & -0.030 \\
\addlinespace
\multicolumn{3}{l}{\textit{GF-PR}} \\
In-approve degree & -0.051 & -0.042 \\
In-approve value (sum) & -0.050 & -0.041 \\

\bottomrule
\end{tabular}
}
\end{table}


EndorseRank는 두 allowance 대리 지표와 신뢰할 만한 일차 구성 개념 정합을 보인다(인승인 차수 $\tau=0.356$; 인승인 가치 $\tau=0.354$; 계열 평균 $0.355$). 전송에서의 AWP와 마찬가지로 이는 일차 구성 개념 검사이다. $G_{\text{approve}}$의 PageRank는 지출자 측 승인 중심성과 상관해야 한다. 따라서 allowance 계열 결과는 부분적으로 동일 구성 개념 검사이며, 독립된 외부 기준이 아니다. 과학적으로 정보적인 칸은 AWP의 영에 가까운 allowance $\tau$로, 전송 중심성이 승인을 복원하지 않음을 보여 준다. EndorseRank의 $\tau \approx 0.355$는 주어진 그래프에서 복원 가능한 그래프 신호의 증거이다. 차수와 가치 상관의 거의 동등함은, 이 코호트에서 많은 소유자에게 승인받는 것과 큰 집계 가치로 승인받는 것이 최신 allowance 의미 아래 유사한 순위 정보를 전달함을 시사한다.

AWP의 allowance 정합은 영에 가깝고 약간 음이다(두 대리 지표 모두 $\tau=-0.030$). 전송 중심성은 지출 승인을 받는 지갑을 식별하지 않는다. 이 분리는 연구 기여의 중심에 있다. 전송 기반 평판이 보증의 충분한 대리 지표라면 AWP는 적어도 중간 정도 allowance 정합을 보여야 한다. 그렇지 않다. 명시적 지출 승인---예를 들어 타인을 대신하여 토큰을 이동하도록 신뢰받는 주소---을 중시하는 프로토콜은, 이 증거에서 AWP를 allowance 인식 방법으로 대체할 수 없다.

allowance 결과는 동일 창 기여의 실증적 핵심이다. 전송 그래프는 지출 승인을 인코딩하지 않는다. 런타임 결과와 결합하면, allowance 계열은 주장~C4를 지지한다. 더 싼 방법이 승인 검사와 정합하는 유일한 사회적 방법이기도 하다. 제~\ref{sec:gf-pr-ceiling}절의 홀드아웃은 그 신호가 점수 창 이후에도 지속되는지를 검정한다.

\dissertationSubsubheading[sub:inverse-risk-family]{역위험 계열 (주장하지 않음)}

표~\ref{tab:alignment-inverse-risk}는 산출물 표이며 본문에서 해석하지 않는다.

% Auto-generated evaluation table fragment — regenerate from the evaluation pipeline; do not edit by hand.

\begin{table}[htbp]
\centering
\caption{Domain alignment with inverse-risk proxies from GMX realized PnL (higher = lower loss).}
\label{tab:alignment-inverse-risk}
\fitwidth{%
\begin{tabular}{lcc}
\toprule
Proxy & Spearman $\rho$ & Kendall $\tau$ \\
\midrule
\multicolumn{3}{l}{\textit{EndorseRank}} \\
Loss avoidance ($-\sum\min(\text{PnL},0)$) & 0.097 & 0.079 \\
Non-loss close rate & -0.037 & -0.031 \\
Worst-close PnL score ($-\min\text{PnL}$) & 0.070 & 0.057 \\
\addlinespace
\multicolumn{3}{l}{\textit{AWP}} \\
Loss avoidance ($-\sum\min(\text{PnL},0)$) & 0.243 & 0.165 \\
Non-loss close rate & 0.028 & 0.019 \\
Worst-close PnL score ($-\min\text{PnL}$) & 0.204 & 0.137 \\
\addlinespace
\multicolumn{3}{l}{\textit{GF-PR}} \\
Loss avoidance ($-\sum\min(\text{PnL},0)$) & 0.675 & 0.545 \\
Non-loss close rate & 0.150 & 0.104 \\
Worst-close PnL score ($-\min\text{PnL}$) & 0.618 & 0.478 \\

\bottomrule
\end{tabular}
}
\end{table}


\dissertationSubsubheading[sub:sybil-family]{Sybil 조정 안정성 계열}

표~\ref{tab:alignment-sybil-stability}는 인바운드 상대방 비율, 전송 재직 기간, 활성 월과의 정합을 보고한다.

% Auto-generated evaluation table fragment — regenerate from the evaluation pipeline; do not edit by hand.

\begin{table}[htbp]
\centering
\caption{Domain alignment with Sybil-adjusted stability proxies from transfer activity.}
\label{tab:alignment-sybil-stability}
\fitwidth{%
\begin{tabular}{lcc}
\toprule
Proxy & Spearman $\rho$ & Kendall $\tau$ \\
\midrule
\multicolumn{3}{l}{\textit{EndorseRank}} \\
Inbound counterparty ratio & 0.093 & 0.073 \\
Transfer tenure (days) & 0.366 & 0.297 \\
Active months & 0.381 & 0.330 \\
\addlinespace
\multicolumn{3}{l}{\textit{AWP}} \\
Inbound counterparty ratio & -0.063 & -0.061 \\
Transfer tenure (days) & 0.602 & 0.436 \\
Active months & 0.622 & 0.482 \\
\addlinespace
\multicolumn{3}{l}{\textit{GF-PR}} \\
Inbound counterparty ratio & -0.339 & -0.235 \\
Transfer tenure (days) & 0.359 & 0.244 \\
Active months & 0.372 & 0.270 \\

\bottomrule
\end{tabular}
}
\end{table}


AWP는 재직 기간과 활성 월에서 EndorseRank를 상회하며, 이는 시간 감쇠 전송 간선에서 예상된다. 인바운드 상대방 비율은 두 방법 모두 약하다. 이 계열은 강건성 노트이며 헤드라인 주장이 아니다.

\dissertationSubsubheading[sub:trading-success-family]{거래 성공 계열 (주장하지 않음)}

표~\ref{tab:alignment-gmx}는 산출물 표이며 본문에서 해석하지 않는다.

% Real BigQuery data -- regenerate with:
%   A-Skill-Programs/margin_rank/scripts/run_dissertation_eval.py --real --export-latex
% Generated by export_latex_results.py -- do not edit by hand unless re-export fails
\begin{table}[htbp]
\centering
\caption[Domain alignment with GMX V2 margin-trading success proxies on Arbitrum One.]{Domain alignment with GMX V2 margin-trading success proxies on Arbitrum One. CI = 95\% percentile bootstrap over wallets (400 paired resamples, seed 42).}
\label{tab:alignment-gmx}
\fitwidth{%
\begin{tabular}{lccc}
\toprule
Proxy & Spearman $\rho$ & Kendall $\tau$ & 95\% CI of $\tau$ \\
\midrule
\multicolumn{4}{l}{\textit{EndorseRank}} \\
Close-success count & 0.180 & 0.149 & [0.130, 0.171] \\
Realized gain proxy & 0.132 & 0.107 & [0.086, 0.127] \\
Close success rate & 0.041 & 0.033 & [0.011, 0.055] \\
\addlinespace
\multicolumn{4}{l}{\textit{AWP}} \\
Close-success count & 0.319 & 0.224 & [0.206, 0.241] \\
Realized gain proxy & 0.349 & 0.239 & [0.223, 0.254] \\
Close success rate & 0.110 & 0.075 & [0.055, 0.092] \\

\bottomrule
\end{tabular}
}
\end{table}


\dissertationSubheading[sec:rank-divergence]{순위 분기 (H1 / RQ1)}

가설 H1과 연구 질문 RQ1은 동일 표본에서 EndorseRank와 AWP가 의미 있게 다른 지갑 순서를 산출하는지를 묻는다. 두 점수 벡터 간 방법 간 Kendall $\tau$는 매칭 코호트($n=5{,}521$)에서 $0.316$이다. 순위는 관련되나 상호 교환과는 거리가 멀다. 경제적으로 활성인 지갑이 전송 시장과 승인 시장 모두에 참여하는 경향이 있어 상관은 존재하지만, allowance 그래프와 전송 그래프는 서로 다른 관계를 강조하므로 많은 쌍별 순서가 불일치한다.

그림~\ref{fig:score-distributions}는 EndorseRank와 AWP 점수의 주변 분포를 보여 준다. 둘 다 희소 방향 그래프의 PageRank에서 전형적이듯 우측 왜도를 보이나, 형상은 동일하지 않다. AWP의 질량은 더 밀집한 전송 연산자를, EndorseRank의 질량은 더 희소한 보증 연산자를 반영한다. 왜도만으로 분기가 증명되지는 않는다---두 방법이 다른 점수 척도로 동일 순서를 공유할 수 있다---그러나 순위--순위 관계를 직접 검토할 동기를 제공한다.

\begin{figure}[htbp]
\centering
\includegraphics[width=0.9\textwidth]{\figpath{score-distributions.pdf}}
\caption{Marginal distributions of EndorseRank and AWP scores on the matched cohort ($n=5{,}521$), shown on a $\log_{10}$ scale for wallets with positive score (EndorseRank excludes $966$ zero-score wallets; AWP excludes $381$). Both are right-skewed PageRank distributions; differences in shape reflect the distinct endorsement and transfer operators.}
\label{fig:score-distributions}
\end{figure}

그림~\ref{fig:rank-scatter}는 EndorseRank 순위를 AWP 순위에 대해 그린다. 대각선 위의 점은 동일 순서를 나타낼 것이고, 관측된 구름은 상당한 비대각 질량을 보여 준다. 전송 중심성에서 높게 순위화되는 지갑이 allowance 보증에서 높게 순위화되는 지갑과 신뢰할 만하게 같지 않으며, 그 반대도 성립한다. 방법 간 $\tau=0.316$은 이 산점을 요약한다. 일치는 우연을 초과하지만, 한 순위를 다른 순위로 대체하는 프로토콜은 지갑 쌍의 큰 비율을 오순위화할 것이다.

\begin{figure}[htbp]
\centering
\includegraphics[width=0.9\textwidth]{\figpath{rank-scatter.pdf}}
\caption{Rank--rank scatter of EndorseRank versus AWP on the matched cohort. Inter-method Kendall $\tau=0.316$ indicates partial overlap with substantial pairwise disagreement, supporting H1/RQ1 (construct-distinct orderings).}
\label{fig:rank-scatter}
\end{figure}

순위 분기는 구성 개념별 정합의 구조적 대응물이다. 두 방법이 거의 동일한 순서를 산출했다면, 표~\ref{tab:alignment-three-methods}의 계열 수준 승자를 구성 개념 효과로 해석하기 어려울 것이다. 공유 순위 주변의 잡음일 수 있다. 관측된 $\tau=0.316$과 EndorseRank의 allowance 우세, AWP의 전송 우세는 방법을 관련되나 구별되는 평판 도구로 읽도록 지지한다. H1은 지지된다. EndorseRank와 AWP는 매칭 코호트에서 부분적으로 중첩되면서도 구성 개념상 구별되는 순서를 산출한다.

$\tau=0.316$의 운영적 읽기는, 임의로 고른 지갑 쌍에서 EndorseRank와 AWP가 어느 지갑이 더 높은지에 합의할 확률이 우연보다 중간 정도로만 높다는 것이다. 현재 전송 기반 점수를 쓰는 프로토콜은 allowance 기반 점수로의 전환---또는 둘의 혼합---이 기존 지갑 순서를 보존한다고 가정할 수 없다. 이전은 쌍의 상당한 비율을 재순위화할 것이며, 가장 큰 불일치는 전송 허브이지만 지출자로 거의 보증되지 않는 지갑, 또는 전송 발자국이 제한된 보증 지출자에 집중될 가능성이 크다. 그림~\ref{fig:rank-scatter}는 그 해석과 일치한다. 구름은 양의 기울기를 가지나 넓다.

분기 결과는 또한 ``온체인 평판''을 단일 대상으로 하는 수사를 규율한다. 이 코호트에서 승인으로부터의 평판과 전송으로부터의 평판은, 상호 교환으로 취급하면 증거를 왜곡할 만큼 실증적으로 구별된다. 제~\ref{ch:discussion}장은 실무적 함의로 돌아간다. 방법 선택은 임의의 PageRank 변형에 대한 일반적 선호가 아니라, 프로토콜이 필요로 하는 신뢰 구성 개념을 따라야 한다.

\dissertationSubheading[sec:gf-pr-ceiling]{미래 신규 승인의 시간 홀드아웃 (RQ4)}

점수는 2026년 2월 28일에 고정한다. 라벨은 2026년 3월 1일부터 2026년 5월 31일까지 처음 관측된 신규 소유자--지출자 승인 쌍이며, 지출자 측에서 센다. 주장 코호트는 t1 인바운드 지출자($n=1{,}335$)이며, 그중 $222$개가 적어도 하나의 신규 t2 쌍을 받았다.

미래 신규 승인자와의 EndorseRank Kendall $\tau$는 $0.479$이다(부트스트랩 95\% 신뢰구간 $0.426$--$0.529$; $200$회 재표본). AWP는 $0.044$이다(신뢰구간 $-0.011$--$0.095$). 결과는 미래 시제의 동일 보증 구성 개념이다. 라벨이 점수화된 이벤트로부터 구축되지 않으므로, 동일 창 allowance $\tau=0.355$보다 창 밖 증거로서 더 강하다. 거래 이익, 청산(liquidation), 악성 지출자의 예측이 아니다.

이 라벨의 GMX 매칭 교집합은 희소하며($12$건) 주장으로 사용하지 않는다. 철회, drain 휴리스틱, Aave 차입자 라벨은 계산되었으나 지지로 보고하지 않는다. 철회와 drain은 보증 읽기에 부호가 틀렸고, 이 지출자 집합의 Aave 이벤트는 비어 있거나 구성 개념이 틀렸다.

\dissertationSubheading[sec:empirical-claims]{주요 주장 (C1--C4)}

제4장 실증 결과는 표~\ref{tab:hrq-claim-map}에 정합하는 다음 주장을 지지한다.

\begin{enumerate}
\item[\textbf{C1.}] 동일 $n$에서 EndorseRank 런타임이 AWP 런타임보다 낮다. 매칭 코호트($n=5{,}521$)에서 EndorseRank는 PageRank를 $2.105$\,s에 완료하고 AWP는 $18.711$\,s에 완료하며(가속 $\approx 8.9\times$), 간선은 $14{,}727$ 대 $137{,}087$개이다(표~\ref{tab:benchmark-runtime}). 확장 지갑 풀에서 가속은 $n=10{,}000$에서 $\approx 10.4\times$($1.18$\,s 대 $12.28$\,s), $N=31{,}612$에서 $\approx 9.4\times$($2.29$\,s 대 $21.40$\,s)로 유지된다(표~\ref{tab:benchmark-scaling}).

\item[\textbf{C2.}] EndorseRank와 AWP는 동일 창 국소 통계에서 의도된 구성 개념을 복원한다. EndorseRank allowance 평균 $\tau=0.355$; AWP 전송 평균 $\tau=0.534$(표~\ref{tab:alignment-transfer}--\ref{tab:alignment-allowance}). 이들은 구성 개념 검사이며 거래 기술 검사가 아니다.

\item[\textbf{C3.}] \textbf{구성 개념별 정합:} EndorseRank는 allowance에서 AWP를 상회하고($\tau=0.355$ 대 $-0.030$), AWP는 전송에서 EndorseRank를 상회한다($\tau=0.534$ 대 $0.333$). 방법 간 순위 일치는 중간 정도에 불과하다($\tau=0.316$)(제~\ref{sec:rank-divergence}절).

\item[\textbf{C4.}] \textbf{효율--정합 상충:} EndorseRank는 매칭 코호트에서 AWP보다 PageRank 런타임이 약 $8.9\times$ 낮고, 확장 풀 스케일링에서 약 $9$--$10\times$ 낮으면서, 승인 검사와 정합하는 유일한 사회적 방법으로 남는다. 지출자 홀드아웃($n=1{,}335$)에서 미래 신규 승인자와의 EndorseRank $\tau=0.479$(신뢰구간 $0.426$--$0.529$) 대 AWP $\tau=0.044$이다(제~\ref{sec:gf-pr-ceiling}절).
\end{enumerate}

RQ1은 방법 간 $\tau=0.316$으로 다룬다. RQ2는 allowance 대 전송 검사로 다룬다. RQ3는 런타임 표로 다룬다. RQ4는 거래 성공 계열이 아니라 지출자 홀드아웃으로 다룬다.

장이 \emph{주장하지 않는} 것 또한 동등하게 중요하다. 어느 점수도 거래 이익, 청산(liquidation), 악성 지출자를 예측한다고 주장하지 않는다. AWP가 구식이 되었다고 주장하지 않는다. AWP는 전송 허브 탐지에서 더 강한 방법으로 남는다. 기여는 구성 개념 인식 비교이다. EndorseRank는 더 싸고, allowance 검사와 정합하며, AWP와 중간 정도로 분기하고, 지출자 홀드아웃에서 이후의 신규 승인을 예측한다.

제~\ref{ch:discussion}장은 이들 발견을 RQ1--RQ4에 비추어 해석하고, 함의, 타당성에 대한 위협, 한계를 논의한다.
